\chapter{Services de positions et de P\&L}\label{ch:pl:position-and-pnl-services}

Trente-trois minutes après le début de la session, la session de saisie des ordres a été coupée pendant quatre secondes. Quatre exécutions ont eu lieu pendant la coupure, et
à son rétablissement, les rapports d'exécution de la stratégie la faisaient apparaître 500 actions moins vendeuse que ne le faisait la drop copy de la plateforme. Le \omterm{def:pl:position-and-pnl-services:service}{service de positions} devait
trancher entre les deux. La réponse est qu'il ne faut croire ni l'une ni l'autre tant qu'elles ne concordent pas~: compter chaque exécution signalée par une source, une seule fois, et dire que la
position est en attente tant que les sources ne concordent pas. Tous les chiffres aval --- risque (chapitre~18), limites, P\&L, les livres qui se clôturent la nuit ---
partent des positions, et les positions partent d'exécutions qui arrivent en double, en retard, dans le désordre, et qui sont parfois annulées.

\section{Positions à partir des exécutions}

\begin{definition}[Service de positions]\label{def:pl:position-and-pnl-services:service}
Un \emph{service de positions}\index{service de positions} (position service) tient à jour les positions et le P\&L de la société par compte et par instrument à partir de ses exécutions,
fusionnées à partir de chaque source qui les signale, et les sert à tout instant, en date de négociation ou en date de règlement, avec un statut qui indique
si ses sources concordent.
\end{definition}

\begin{definition}[Identifiant d'exécution]\label{def:pl:position-and-pnl-services:execid}
Un \emph{identifiant d'exécution}\index{identifiant d'exécution} (execution identifier) est l'identifiant unique attribué par la plateforme à une exécution (dans le simulateur, son numéro
d'appariement avec l'identifiant client de l'ordre)~; c'est lui qui permet de reconnaître sans ambiguïté une exécution signalée par deux sources, ou deux fois par une même source.
\end{definition}

Le service du chapitre (\texttt{firm.posservice}) stocke les exécutions indexées par cet identifiant et ne stocke jamais de position~: chaque position est
calculée en rejouant les exécutions dans l'ordre chronologique au moyen du gestionnaire exact du livre~1 (\texttt{firm.pnl}), où la quantité et le cash sont des
entiers, de sorte que le cash plus la quantité multipliée par la valorisation, moins les frais, soit exact à la dix-millième de dollar près. Stocker les événements plutôt que leur
somme rend l'ordre d'arrivée sans importance, rend les doublons inoffensifs (traitement idempotent, livre~13, chapitre~24) et rend toute position passée
recalculable.

\omcode{code/firm/posservice/firm_posservice.py}{59}{71}{Exécutions venant de n'importe quelle source, comptées une seule fois~; annulations et corrections conservées comme événements,
non appliquées en modifiant le passé.}\label{lst:pl:position-and-pnl-services:events}

\omcode{code/firm/posservice/firm_posservice.py}{74}{97}{Une position est un rejeu~: les exécutions en vigueur à un instant donné (annulations et corrections
appliquées si elles sont connues à cet instant), dans l'ordre chronologique, à travers le gestionnaire exact.}\label{lst:pl:position-and-pnl-services:position}

\begin{omfigure}
\begin{tikzpicture}[font=\scriptsize,box/.style={draw,rounded corners=2pt,align=center,minimum height=0.8cm,inner sep=3pt},>=Stealth]
\node[box,fill=omMeth!12,minimum width=2.8cm] (s) at (0,1.3) {session de saisie\\ des ordres~: rapports\\ d'exécution};
\node[box,fill=omMeth!12,minimum width=2.8cm] (d) at (0,0) {drop copy~: chaque\\ exécution, avec\\ un léger retard};
\node[box,fill=omMeth!12,minimum width=2.8cm] (b) at (0,-1.3) {avis de la plateforme~:\\ annulations,\\ corrections};
\node[box,fill=omDef!15,minimum width=3.4cm,minimum height=2.6cm] (p) at (4.6,0) {service de positions\\ exécutions indexées par\\ identifiant d'exécution\\ statut~: concordant\\ ou en attente};
\node[box,fill=gray!10,minimum width=3.4cm] (v) at (9.2,1.2) {vues~: date de négociation,\\ date de règlement~; valorisations~:\\ intrajournalières, de fin de journée};
\node[box,fill=gray!10,minimum width=3.4cm] (r) at (9.2,-1.2) {rapprochement avec\\ le relevé de compensation};
\foreach \x in {s,d,b} {\draw[->] (\x.east) -- (p.west|-\x.east);}
\draw[->] (p.east) -- (v.west); \draw[->] (p.east) -- (r.west);
\end{tikzpicture}

\omcaption{Le \omterm{def:pl:position-and-pnl-services:service}{service de positions}. Les exécutions arrivent de la session et de la drop copy, les annulations et corrections de la plateforme~; le service
conserve les exécutions, non les positions, et calcule chaque vue à partir d'elles~; le relevé de la chambre de compensation est le contrôle en
bout de chaîne.}\label{fig:pl:position-and-pnl-services:service}
\end{omfigure}

\section{Deux flux d'une même réalité~: sessions et drop copies}

La session de saisie des ordres signale les exécutions des ordres qu'elle a envoyés~; la drop copy (livre~13, chapitre~21) envoie à la société une copie de chaque
exécution de ses sessions sur une connexion distincte, précisément au moment où la première tombe. Une session qui se coupe puis se reconnecte en demandant
uniquement les nouveaux messages ne reçoit jamais les rapports générés pendant la coupure~; une drop copy qui accuse un léger retard (50~millisecondes ici) les a. Le service prend
les deux, compte chaque exécution une seule fois et compare les deux ensembles~: lorsqu'ils diffèrent, la position qu'il sert est l'union,
et son statut est en attente.

\omcode{code/firm/posservice/firm_posservice.py}{99}{107}{Le statut~: concordant lorsque chaque source a signalé les mêmes exécutions, en attente
sinon, avec les exécutions qui manquent à chaque source.}\label{lst:pl:position-and-pnl-services:status}

\section{Annulations, corrections et exécutions tardives}

\begin{definition}[Annulation de transaction]\label{def:pl:position-and-pnl-services:bust}
Une \emph{annulation de transaction}\index{annulation de transaction} (trade bust) est l'annulation d'une exécution par la plateforme après qu'elle a eu lieu --- par exemple au titre de
ses règles sur les transactions manifestement erronées (livre~1, chapitre~31) --- de sorte que, pour le relevé de la chambre de compensation, la transaction n'a jamais eu lieu.
\end{definition}

Une annulation arrive quelques minutes ou quelques heures après l'exécution~; une correction modifie le prix ou la quantité d'une exécution. Le service applique les deux
comme des événements horodatés, de sorte que la position à un instant quelconque puisse être calculée telle qu'elle était connue alors (l'annulation pas encore reçue)
ou telle qu'elle est connue maintenant (la transaction n'a jamais eu lieu). Une exécution tardive --- un rapport d'exécution qui arrive après que la position a été lue --- n'est qu'un
événement en retard~: le rejeu le remet à sa place.

\section{P\&L intrajournalier et valorisations de fin de journée}

\begin{definition}[P\&L intrajournalier, valorisation de fin de journée]\label{def:pl:position-and-pnl-services:marks}
Le \emph{P\&L intrajournalier}\index{P\&L intrajournalier} (intraday P\&L) est le P\&L des positions valorisées aux cours de marché courants au cours de la journée, pour le suivi et les limites.
La \emph{valorisation de fin de journée}\index{valorisation de fin de journée} (end-of-day mark) est le prix auquel les positions sont valorisées pour le P\&L officiel et les livres de la journée, généralement
le cours de clôture officiel, fixé par la politique de valorisation plutôt que par la dernière cotation que la stratégie a eu l'occasion de voir.
\end{definition}

Le chapitre valorise la journée au mid du flux et la clôture au dernier échange de la plateforme, qui tient lieu de cours de clôture officiel. Le P\&L est
exact~: \texttt{firm.pnl} conserve le cash et la quantité sous forme d'entiers, et les parts réalisée et latente selon le coût moyen (livre~1, chapitre~7)
s'additionnent, frais compris, pour donner le total à la valorisation.

\section{Début de journée et livres qui se clôturent}

\begin{definition}[Position de début de journée, position en date de négociation, position en date de règlement]\label{def:pl:position-and-pnl-services:sod}
La \emph{position de début de journée}\index{position de début de journée} (start-of-day position) est la position avec laquelle une journée commence~: la position de fin de journée de la veille,
rapprochée. La \emph{position en date de négociation}\index{position en date de négociation} (trade-date position) compte chaque exécution à partir du moment où elle se négocie~; la
\emph{position en date de règlement}\index{position en date de règlement} (settlement-date position) ne la compte qu'à partir de sa date de règlement (un jour ouvré plus tard pour les actions américaines,
livre~1, chapitre~5), ce que reflètent la conservation des titres et le cash.
\end{definition}

À la fin de la journée, le service enregistre chaque position avec sa \omterm{def:pl:position-and-pnl-services:marks}{valorisation de fin de journée} et son total~; la \omterm{def:pl:position-and-pnl-services:sod}{position de début de journée} du lendemain est cet
enregistrement. Le relevé de la chambre de compensation --- les exécutions qu'elle a compensées, les annulations retirées --- est le contrôle~: le service rapproche ses positions
de ce relevé et signale chaque écart de rapprochement (break) avant que la journée suivante ne commence.

\section{Quatre secondes sans session}

L'expérience du chapitre fait tourner une stratégie de cotation (500 actions de chaque côté, une limite d'inventaire de 5\,000 actions, ordres laissés en attente si la session
se coupe) pendant une heure sur dix sessions \texttt{firm.tape} du simulateur du livre~10, chacune avec une session de drop copy. À une minute tirée au hasard, la session de saisie des ordres
est coupée pendant quatre secondes~; le client se reconnecte en ne demandant que les nouveaux messages, et l'équipe des opérations demande un rejeu des rapports manqués une
minute plus tard. À la minute~50, la plateforme annule l'exécution la plus proche de la minute~30. Trois services sont comparés au relevé de la chambre de compensation~: rapports de session
seuls~; session et drop copy fusionnées~; fusionnées avec l'annulation appliquée.

\omcode{code/platforms/17-position-and-pnl-services/python/pl_posservice.py}{72}{91}{Les trois services alimentés dans l'ordre d'arrivée~: les rapports de la session
(les rapports perdus seulement au rejeu), la drop copy avec 50~millisecondes de retard, et l'avis d'annulation.}\label{lst:pl:position-and-pnl-services:feed}

\begin{omfigure}
\begin{tikzpicture}
\begin{axis}[width=12cm,height=5.4cm,xlabel={minutes depuis le début de la session},ylabel={position moins relevé (actions)},xmin=25,xmax=60,
  ymin=-50,ymax=700,grid=major,grid style={gray!20},tick label style={font=\small},label style={font=\small},table/col sep=comma]
\addplot[color=omThm,very thick,const plot] table[x=t_min,y expr=\thisrow{session}-\thisrow{truth}]
  {figdata/platforms/17-position-and-pnl-services/timeline.csv};
\addplot[color=omDef,thick,const plot] table[x=t_min,y expr=\thisrow{merged}-\thisrow{truth}]
  {figdata/platforms/17-position-and-pnl-services/timeline.csv};
\addplot[color=black,thick,dashed,const plot] table[x=t_min,y expr=\thisrow{merged_bust}-\thisrow{truth}]
  {figdata/platforms/17-position-and-pnl-services/timeline.csv};
\node[omThm,font=\scriptsize,anchor=north west] at (axis cs:34,585) {session seule~: quatre exécutions perdues à la minute 33};
\node[omDef,font=\scriptsize,anchor=south west] at (axis cs:50.5,120) {avec drop copy};
\node[font=\scriptsize,anchor=east] at (axis cs:49.5,45) {avec drop copy et annulation};
\end{axis}
\end{tikzpicture}

\omcaption{La position de chaque service moins celle de l'état de compensation, pendant l'heure de la graine~3 (échantillonnée toutes les cinq secondes). L'exécution annulée
de la minute~30 écarte chaque service de 100 actions de l'état~; la panne de la minute~33 ajoute 500 au service de session seule~; le service avec le drop copy et l'annulation
concorde avec l'état à partir de l'avis d'annulation de la minute~50. Données~:
\texttt{fig\_posservice.py}.}\label{fig:pl:position-and-pnl-services:timeline}
\end{omfigure}

\begin{example}[Graine 3]\label{ex:pl:position-and-pnl-services:seed3}
La session tombe à la minute~33 et quatre exécutions ont lieu en quatre secondes~; le service de session seule est écarté de 500 actions à partir de là. Le
service fusionné les reçoit 50~millisecondes après qu'elles ont eu lieu, par le drop copy, et signale sa position comme en attente pendant 62~secondes, jusqu'à ce que le
rejeu remette la session en accord. Les deux portent l'exécution de la minute~30, que la plateforme annule à la minute~50~; seul le service qui
applique l'avis concorde avec l'état de compensation à la clôture. À la clôture, l'état affiche $-3\,300$ actions et $-\$3\,493$~; le
service de session seule est écarté de 600 actions et de 252~\$, le service fusionné de 100 actions et de 47~\$, le service fusionné avec l'annulation est exact.
\end{example}

\begin{table}[ht]
\centering\small
\begin{tabular}{lrr}
\toprule
service & erreur moyenne $|$de position$|$ (actions) & erreur moyenne $|$de P\&L$|$ (\$) \\
\midrule
session seule & 290 & 52.4 \\
avec le drop copy & 220 & 28.4 \\
avec le drop copy et l'annulation & 0 & 0.0 \\
\bottomrule
\end{tabular}
\caption{Erreurs à la clôture par rapport à l'état de compensation, moyennées sur dix sessions d'une heure, chacune avec une panne de quatre secondes à une
minute tirée au sort et une exécution annulée vingt minutes après sa survenue.}\label{tab:pl:position-and-pnl-services:close}
\end{table}

Sur les dix sessions (\cref{tab:pl:position-and-pnl-services:close}), la panne coûte peu exécution par exécution --- elle touche
des exécutions dans quatre sessions (une, quatre, trois et une), car en quatre secondes une stratégie de cotation négocie rarement --- et l'annulation coûte
quelque chose chaque jour. Sans le drop copy, le service termine la journée en moyenne à 290 actions et à 52~\$ d'écart~; avec lui, à 220 actions et 28~\$ d'écart, dont
la totalité provient de l'annulation~; avec l'annulation appliquée, il est exact. Dans les sessions où la panne a fait perdre des exécutions, le service fusionné a signalé une position en attente pendant 60 à
63~secondes, jusqu'au rejeu~; sans rejeu, il l'aurait signalée en attente jusqu'à ce que quelqu'un regarde.

\begin{omfigure}
\begin{tikzpicture}
\begin{axis}[width=12cm,height=5.4cm,ybar,bar width=5pt,ylabel={erreur de position à la clôture (actions)},
  xlabel={session (graine)},xtick={1,2,3,4,5,6,7,8,9,10},ymin=-700,ymax=700,grid=major,grid style={gray!20},tick label style={font=\small},
  label style={font=\small},table/col sep=comma,legend style={font=\footnotesize,at={(0.5,1.03)},anchor=south,legend columns=3}]
\addplot[fill=omThm!45,draw=omThm] table[x=seed,y=session_pos_err]{figdata/platforms/17-position-and-pnl-services/at_close.csv};
\addlegendentry{session seule}
\addplot[fill=omDef!45,draw=omDef] table[x=seed,y=merged_pos_err]{figdata/platforms/17-position-and-pnl-services/at_close.csv};
\addlegendentry{avec drop copy}
\addplot[fill=black!30,draw=black] table[x=seed,y=merged_bust_pos_err]{figdata/platforms/17-position-and-pnl-services/at_close.csv};
\addlegendentry{avec drop copy et annulation}
\end{axis}
\end{tikzpicture}

\omcaption{Position de chaque service moins celle de l'état de compensation à la clôture, session par session. Sans le drop copy, l'erreur est la
somme des exécutions perdues et de l'exécution annulée~; avec lui, la seule exécution annulée~; avec l'annulation appliquée, rien (la troisième barre est nulle partout).
Données~: \texttt{fig\_posservice.py}.}\label{fig:pl:position-and-pnl-services:close}
\end{omfigure}

\begin{remark}[Ce qu'apporte l'attente]\label{rem:pl:position-and-pnl-services:pending}
Un statut en attente transforme une erreur silencieuse en erreur visible. La position du service fusionné était juste pendant la panne, mais seul le
désaccord entre ses sources a signalé que quelque chose était arrivé à la session~; le service de session seule était faux et sûr de lui. Le risque
(chapitre~18) doit traiter une position en attente comme l'union de ce que rapporte n'importe quelle source --- l'exposition la plus grande --- jusqu'à ce qu'elle concorde.
\end{remark}

\section{Tutoriel~: une journée, trois sources, un état}\label{tut:pl:position-and-pnl-services:tut}

\begin{tutorial}
\textbf{Objectif.} Construire les positions à partir de deux flux d'exécutions et d'une annulation, observer leur désaccord puis leur accord, et les rapprocher de l'état de compensation.
\textbf{État final~:} \cref{fig:pl:position-and-pnl-services:timeline} et \cref{tab:pl:position-and-pnl-services:close}.
\begin{tutsteps}
\item \textbf{Une journée}~: \texttt{pl\_posservice.day(seed)} lance le simulateur avec une session de drop copy et une panne programmée.
\item \textbf{Services}~: \texttt{services(d)}~; \texttt{status} et \texttt{missing} pendant la panne.
\item \textbf{L'état}~: \texttt{truth(d)}~; \texttt{reconcile(service, statement)}.
\item \textbf{La clôture}~: \texttt{at\_close()} sur dix graines~; \texttt{end\_of\_day} et \texttt{start\_of\_day} pour le roll.
\end{tutsteps}
\textbf{Pour aller plus loin.} Annuler les ordres à la déconnexion et voir ce que coûte alors la panne~; retarder le drop copy d'une seconde et observer les
fenêtres d'attente s'allonger.
\end{tutorial}

\section{Construction~: le service de positions}\label{bld:pl:position-and-pnl-services:posservice}

\begin{build}
\bfield{Rôle} Une source unique de positions et de P\&L pour le risque, les limites, le reporting et les livres, correcte face aux rapports en double, en retard ou manquants,
aux annulations et aux corrections, et honnête lorsqu'elle n'est pas sûre.
\bfield{Interface} \texttt{Execution}~; \texttt{PositionService(sources)} avec \texttt{on\_execution}, \texttt{on\_bust},
\texttt{on\_correction}, \texttt{position(account, symbol, as\_of, view, date)}, \texttt{status}, \texttt{missing}, \texttt{pnl},
\texttt{end\_of\_day}, \texttt{start\_of\_day}~; \texttt{reconcile}.
\bfield{Règles} Les exécutions sont stockées, les positions calculées~; un \omterm{def:pl:position-and-pnl-services:execid}{identifiant d'exécution} ne compte qu'une fois~; les annulations et les corrections sont des événements datés~;
quantités entières et espèces (\texttt{firm.pnl})~; en attente dès que les sources se contredisent.
\bfield{Tests d'acceptation} \texttt{code/firm/posservice/tests/}~: fusion et statut idempotents~; P\&L exact et indépendant de l'ordre d'arrivée~;
annulations et corrections datées~; vues en date de négociation et en date de règlement, roll de fin de journée et écarts de rapprochement.
\bfield{Pour aller plus loin} De nombreux comptes et instruments, avec des allocations à partir d'un compte à coût moyen~; opérations sur titres touchant les positions (chapitre~6)~;
les événements du service journalisés et rejoués (chapitre~12).
\end{build}

\begin{omsources}
\item One Quant Book~1, chapitres 5, 7 et 31 (règlement-livraison, positions et P\&L, transactions erronées)~; One Quant Book~13, chapitres 21 et 24 (drop
copies, traitement idempotent)~; One Quant Book~10, chapitre~26 (le drop copy du simulateur et les pannes de session).
\end{omsources}

\section{Exercices}

\begin{exercise}[$\star$]\label{exo:pl:position-and-pnl-services:1}
Une même exécution arrive de la session, puis deux fois du drop copy. Déterminer la position et en expliquer la raison.
\end{exercise}

\begin{exercise}[$\star$]\label{exo:pl:position-and-pnl-services:2}
Un achat de 100 à 100.00~\$ et une vente de 300 à 100.02~\$, valorisés à 100.00~\$. Calculer la position, le P\&L réalisé et le P\&L latent.
\end{exercise}

\begin{exercise}[$\star$]\label{exo:pl:position-and-pnl-services:3}
Expliquer pourquoi une transaction annulée est retirée comme un événement plutôt que supprimée du stockage.
\end{exercise}

\begin{exercise}[$\star\star$]\label{exo:pl:position-and-pnl-services:4}
Pourquoi l'interruption n'a-t-elle fait perdre des exécutions que dans quatre des dix sessions, et qu'est-ce qui changerait cela~?
\end{exercise}

\begin{exercise}[$\star\star$]\label{exo:pl:position-and-pnl-services:5}
Une transaction du vendredi se règle le jour ouvré suivant. Que montrent les vues en date de négociation et en date de règlement le vendredi soir et le lundi~?
\end{exercise}

\begin{exercise}[$\star\star$]\label{exo:pl:position-and-pnl-services:6}
La position du service fusionné était juste pendant l'interruption. Pourquoi son statut en attente reste-t-il utile~?
\end{exercise}

\begin{exercise}[$\star\star\star$]\label{exo:pl:position-and-pnl-services:7}
\emph{Programmation.} Exécuter \texttt{at\_close} avec la stratégie qui annule ses ordres à la déconnexion. Que deviennent les exécutions perdues pendant l'interruption~?
\end{exercise}

\begin{exercise}[$\star\star\star$]\label{exo:pl:position-and-pnl-services:8}
\emph{Trouver la faille.} «~Notre \omterm{def:pl:position-and-pnl-services:service}{service de positions} ajoute chaque exécution à la position dès son arrivée, et se recale sur le fichier de compensation chaque matin, donc
il est juste au moins une fois par jour.~»
\end{exercise}

\section{Problème~: quatre secondes sans session}

\begin{problem}[{Problème du week-end --- positions issues de deux flux et une annulation de transaction}]\label{pb:pl:position-and-pnl-services:1}
La stratégie du chapitre, dix sessions, une interruption de quatre secondes et une \omterm{def:pl:position-and-pnl-services:bust}{annulation de transaction} (trade bust) par session.

\medskip\noindent\textbf{Partie I --- Le service.}
\begin{enumerate}
\item Que stocke le service, et comment calcule-t-il une position~?
\item Qu'est-ce qui fait que deux rapports désignent la même exécution~?
\item Que dit le statut, et quand~?
\item Comment le service traite-t-il une \omterm{def:pl:position-and-pnl-services:bust}{annulation de transaction} et une correction~?
\item Qu'est-ce qui garantit l'exactitude du P\&L~?
\end{enumerate}

\medskip\noindent\textbf{Partie II --- Les sources.}
\begin{enumerate}[resume]
\item Pourquoi une session perd-elle des rapports lors d'une interruption~?
\item Qu'apporte le drop copy, et à quel coût~?
\item Qu'est-ce qu'un rejeu, et pourquoi le demande-t-on~?
\item Que représente le relevé de compensation~?
\item Quand chaque service concorde-t-il avec le relevé~?
\end{enumerate}

\medskip\noindent\textbf{Partie III --- La journée.}
\begin{enumerate}[resume]
\item Dans combien de sessions l'interruption a-t-elle fait perdre des exécutions, et combien~?
\item Combien de temps le service fusionné est-il resté en attente~?
\item Quelles ont été les erreurs à la clôture dans la graine~3~?
\item Quelles ont été les erreurs moyennes sur les dix sessions~?
\item Quelle erreur le drop copy supprime-t-il, et laquelle non~?
\end{enumerate}

\medskip\noindent\textbf{Partie IV --- Le verdict.}
\begin{enumerate}[resume]
\item Énoncer le \emph{résultat nommé}~: l'erreur de position et de P\&L à la clôture sans le drop copy, avec lui, et avec lui plus l'\omterm{def:pl:position-and-pnl-services:bust}{annulation de transaction},
ainsi que le délai avant que le service concorde avec le relevé de compensation.
\item Que doit faire le risque d'une position en attente~?
\item Que doit-il se passer au début du jour suivant~?
\item Que surveilleriez-vous sur le \omterm{def:pl:position-and-pnl-services:service}{service de positions} lui-même~?
\item En une phrase~: qu'est-ce qui rend une position fiable~?
\end{enumerate}
\end{problem}

\section{Questions d'entretien}

\begin{interviewq}[$\star$ \iqroles{développeur}]\label{iq:pl:position-and-pnl-services:1}
Comment rendez-vous un \omterm{def:pl:position-and-pnl-services:service}{service de positions} robuste aux rapports d'exécution en double~?
\end{interviewq}

\begin{interviewq}[$\star\star$ \iqroles{développeur}]\label{iq:pl:position-and-pnl-services:2}
Votre session de saisie des ordres a été interrompue pendant quelques secondes. Comment savez-vous que votre position est juste~?
\end{interviewq}

\begin{interviewq}[$\star\star$ \iqroles{développeur, chercheur}]\label{iq:pl:position-and-pnl-services:3}
Pourquoi stocker les exécutions plutôt que les positions~?
\end{interviewq}

\begin{interviewq}[$\star\star$ \iqroles{développeur}]\label{iq:pl:position-and-pnl-services:4}
Quelle est la différence entre les \omterm{def:pl:position-and-pnl-services:sod}{positions en date de négociation} et en date de règlement, et qui utilise chacune~?
\end{interviewq}

\begin{interviewq}[$\star\star\star$ \iqroles{développeur}]\label{iq:pl:position-and-pnl-services:5}
Concevez le \omterm{def:pl:position-and-pnl-services:service}{service de positions} et de P\&L d'une société de trading multistratégie.
\end{interviewq}
